% Options for packages loaded elsewhere
\PassOptionsToPackage{unicode}{hyperref}
\PassOptionsToPackage{hyphens}{url}
%
\documentclass[
]{article}
\usepackage{amsmath,amssymb}
\usepackage{iftex}
\ifPDFTeX
  \usepackage[T1]{fontenc}
  \usepackage[utf8]{inputenc}
  \usepackage{textcomp} % provide euro and other symbols
\else % if luatex or xetex
  \usepackage{unicode-math} % this also loads fontspec
  \defaultfontfeatures{Scale=MatchLowercase}
  \defaultfontfeatures[\rmfamily]{Ligatures=TeX,Scale=1}
\fi
\usepackage{lmodern}
\ifPDFTeX\else
  % xetex/luatex font selection
\fi
% Use upquote if available, for straight quotes in verbatim environments
\IfFileExists{upquote.sty}{\usepackage{upquote}}{}
\IfFileExists{microtype.sty}{% use microtype if available
  \usepackage[]{microtype}
  \UseMicrotypeSet[protrusion]{basicmath} % disable protrusion for tt fonts
}{}
\makeatletter
\@ifundefined{KOMAClassName}{% if non-KOMA class
  \IfFileExists{parskip.sty}{%
    \usepackage{parskip}
  }{% else
    \setlength{\parindent}{0pt}
    \setlength{\parskip}{6pt plus 2pt minus 1pt}}
}{% if KOMA class
  \KOMAoptions{parskip=half}}
\makeatother
\usepackage{xcolor}
\usepackage[margin=1in]{geometry}
\usepackage{graphicx}
\makeatletter
\def\maxwidth{\ifdim\Gin@nat@width>\linewidth\linewidth\else\Gin@nat@width\fi}
\def\maxheight{\ifdim\Gin@nat@height>\textheight\textheight\else\Gin@nat@height\fi}
\makeatother
% Scale images if necessary, so that they will not overflow the page
% margins by default, and it is still possible to overwrite the defaults
% using explicit options in \includegraphics[width, height, ...]{}
\setkeys{Gin}{width=\maxwidth,height=\maxheight,keepaspectratio}
% Set default figure placement to htbp
\makeatletter
\def\fps@figure{htbp}
\makeatother
\setlength{\emergencystretch}{3em} % prevent overfull lines
\providecommand{\tightlist}{%
  \setlength{\itemsep}{0pt}\setlength{\parskip}{0pt}}
\setcounter{secnumdepth}{-\maxdimen} % remove section numbering
\ifLuaTeX
  \usepackage{selnolig}  % disable illegal ligatures
\fi
\usepackage{bookmark}
\IfFileExists{xurl.sty}{\usepackage{xurl}}{} % add URL line breaks if available
\urlstyle{same}
\hypersetup{
  hidelinks,
  pdfcreator={LaTeX via pandoc}}

\author{}
\date{\vspace{-2.5em}}

\begin{document}

\section{Question}\label{question}

Un estudiante tiene un vaso de forma cilíndrica. El vaso tiene una base
circular de radio 5 cm y una altura de 5 cm, como se muestra en la
figura.

\begin{center}
\\begin{tikzpicture}[scale=0.7]
  % Dibujar cilindro
  \\draw[thick] (0,0) ellipse (5 and 1);
  \\draw[thick] (0,5) ellipse (5 and 1);
  \\draw[thick] (-5,0) -- (-5,5);
  \\draw[thick] (5,0) -- (5,5);
  \\draw[dashed] (0,0) -- (5,0);
  
  % Etiquetas y dimensiones
  \\draw[<->] (5.5,0) -- (5.5,5) node[midway,right] {5 cm};
  \\draw[<->] (-5,-1.5) -- (5,-1.5) node[midway,below] {10 cm};
  \\node at (2.5,-0.5) {5 cm};
  
  % Texto de la operación
  \\node[align=center] at (5,2.5) {$\\pi \\times 5^2 \\times 5 = 392.7$};
\\end{tikzpicture}
\end{center}

A partir de la información anterior, el estudiante plantea la siguiente
operación:

π × 5² × 5 = 392.7

¿A qué corresponde el resultado de la anterior operación?

\subsection{Answerlist}\label{answerlist}

\begin{itemize}
\tightlist
\item
  Al volumen del vaso.
\item
  Al área de la tapa del vaso.
\item
  Al perímetro de la tapa del vaso.
\item
  Al área lateral del vaso.
\end{itemize}

\section{Solution}\label{solution}

La operación π × 5 ² × 5 = 392.7 corresponde al \textbf{área lateral del
vaso}.

Fórmula del área lateral de un cilindro: A = 2π × r × h

Donde: - r es el radio de la base = 5 cm - h es la altura = 5 cm

Sustituyendo: A = 2π × 5 × 5 = 157.08 cm²

\subsection{Answerlist}\label{answerlist-1}

\begin{itemize}
\tightlist
\item
  Falso
\item
  Falso
\item
  Falso
\item
  Verdadero
\end{itemize}

\section{Meta-information}\label{meta-information}

exname: Volumen de un cilindro extype: schoice exsolution: 0001
exshuffle: TRUE

\end{document}
